\begin{center}\begin{tikzpicture}[font=\small,>=stealth,line join=round]
\draw[gray!40,rounded corners=3pt] (0,0) rectangle (14.2,8.3);\node at (7.1,7.700000000000001) {两条对称轴：反射两次，对应一个平移};
\begin{scope}[shift={(2.330769230769231,3.7966666666666664)},x=1.5384615384615383cm,y=1.5333333333333332cm]\draw[->,gray] (-0.8,0)--(3.1,0) node[right]{$x$};\draw[->,gray] (0,-1.4)--(0,1.6) node[above]{$y$};\draw[blue,thick,solid,domain=-0.5:2.8,samples=180] plot(\x,{cos(deg(pi*\x))});\fill[orange!85!black] (0.25,0.7071067811865476) circle(2.2pt);\fill[orange!85!black] (-0.25,0.7071067811865476) circle(2.2pt);\fill[orange!85!black] (2.25,0.7071067811865476) circle(2.2pt);\draw[violet,dashed] (0,-1.3)--(0,1.3);\draw[violet,dashed] (1,-1.3)--(1,1.3);\node[above,violet] at (1,1.3){$x=1$};\end{scope}\node[anchor=west,align=left,text width=5.0cm] at (8.4,4.0) {图取 $f(x)=\cos(\pi x)$\\[9pt]两轴 $x=0,\ x=1$\\[9pt]横坐标配对\\[9pt]$0.25\to-0.25\to2.25$\\[9pt]最后相差 $2$};
\node[text width=12.6cm,align=center] at (7.1,.65) {先关于0反射，再关于1反射；一般得到周期 $2|b-a|$，未必最小。};
\end{tikzpicture}\end{center}
